\documentclass{../../../unipd-notes}

\unipdsetup{
  course = {Fondamenti di Ingegneria Informatica},
  short-course = {Fondamenti di Ingegneria},
  document-type = {Appunti delle lezioni}
}

\begin{document}
\chapter*{Rappresentazione binaria dell'informazione}

Un calcolatore rappresenta dati e istruzioni mediante sequenze di bit. Con $n$ bit si possono codificare $2^n$ configurazioni distinte; l'interpretazione assegnata a ciascuna configurazione dipende dal tipo del dato e dal formato adottato.

La notazione \emph{complemento a due} consente di rappresentare numeri interi con segno. Il bit più significativo assume peso negativo e l'intervallo rappresentabile è \textbf{$[-2^{n-1},2^{n-1}-1]$}. Identificatori tecnici come \texttt{uint32\_t}, la funzione \texttt{decodeInstruction()} e il percorso \texttt{/dev/ttyUSB0} sono composti con carattere monospaziato.

\section*{Procedura di conversione}
Per convertire un intero negativo in complemento a due si applicano i seguenti passaggi.

\begin{itemize}
  \item fissare la lunghezza della parola;
  \item rappresentare il valore assoluto in binario;
  \item complementare i bit e aggiungere uno.
\end{itemize}

\begin{enumerate}
  \item scrivere $13_{10}=00001101_2$ su otto bit;
  \item invertire i bit ottenendo $11110010_2$;
  \item sommare uno: $-13_{10}=11110011_2$.
\end{enumerate}

\section*{Osservazione}
\begin{displayquote}
L'estensione del segno replica il bit più significativo e conserva il valore dell'intero rappresentato.
\end{displayquote}
\end{document}
